% AI-written, AI cross-reviewed; no human sign-off; not a claim that RH is proved
% Atlas unit D481. Local manuscript v3, item G.501, heading on printed PDF p. 282.
% Formalism.L2.DIV.closed_linkDIVt_of_closed_bridge_kernel
% paper2/decomposition_tables.tex:4152–4157. Verbatim excerpt; SHA-256 d302b9d0773fb0d50c94c5c3ea141c80d976cd12141d857b75b720da4f53fd0d.
% Original TeX uses the macros and cross-reference keys of the source manuscript.

\begin{formalstatement}[D481]\label{decomp:D481}
Suppose a family class has a tail of division-kernel data and a closed uniform bound for those kernels whenever the tuple and discriminant-size estimates hold. Under these size estimates the covered-twist and eta-sum bounds imply the division height, frame and slope bounds on a common tail.
\end{formalstatement}
{\small\raggedright\noindent\textsf{\bfseries Routine assembly}\quad \nolinkurl{Formalism.L2.DIV.closed_linkDIVt_of_closed_bridge_kernel}\par\noindent \nolinkurl{Development/AutomorphicArithmetic/ThetaCorrespondence/ThetaPairLifts/ThirdPlaceSeries/DIVRecBridgeSel.lean:197--205}\par\noindent SHA-256 \nolinkurl{ba4779fe82b80ed4d297ed2593dcebc96ca4e68ce674e24e37713d6f6cb3a4e9}\par\noindent Dominated 5,437\quad Disjoint 5,437\quad terminal.\par}
Applies the single-family bridge-kernel implication to each closed witness and keeps its thresholds.
{\small\raggedright Local proof uses \nolinkurl{Formalism.L2.DIV.linkDIVt_of_bridge_kernel}.\par}
